\documentclass[letterpaper,11pt]{article}
\usepackage[spanish,mexico]{babel}
\usepackage[utf8]{inputenc}
\usepackage[T1]{fontenc}
\usepackage{lmodern,microtype}
\usepackage{graphicx,float,placeins}
\usepackage{booktabs,array,tabularx,longtable}
\usepackage[table]{xcolor}
\usepackage{amsmath,enumitem}
\usepackage{setspace,parskip}


\usepackage[
    left=25mm, 
    right=25mm,
    top=35mm,
    bottom=30mm,
    headheight=77pt, 
    papersize={21.59cm,27.94cm}
]{geometry}


\usepackage{fancyhdr}
\usepackage{caption}
\usepackage{hyperref}
\usepackage[notocbib]{apacite}


\hypersetup{colorlinks=true,linkcolor=blue,citecolor=blue,filecolor=magenta,urlcolor=cyan}
\hypersetup{pdftitle={Práctica 3: Machine Learning II},pdfauthor={Lara Hernandez, Lugo Manzano y Garcia Cortes}}


\newcolumntype{L}[1]{>{\raggedright\arraybackslash}p{#1}}
\newcolumntype{Y}{>{\raggedright\arraybackslash}X}


\captionsetup{font=small,labelfont=bf,skip=7pt}
\setlength{\parskip}{5pt}
\setlength{\emergencystretch}{2em}
\setlist{nosep,leftmargin=*}
\setcounter{tocdepth}{2}
\pagestyle{fancy}
\fancyhf{} 


\fancyhead[L]{\includegraphics[height=1.5cm]{portada_img/der.png}}
\fancyhead[R]{\includegraphics[height=1.5cm]{portada_img/izq.png}}
\fancyfoot[C]{\thepage} 


\renewcommand{\headrulewidth}{0.4pt} 
\renewcommand{\footrulewidth}{0pt}  

\fancypagestyle{indice}{
  \fancyhf{}
  \fancyhead[L]{\includegraphics[height=1.5cm]{portada_img/der.png}}
  \fancyhead[R]{\includegraphics[height=1.5cm]{portada_img/izq.png}}
  \fancyfoot{}
  \renewcommand{\headrulewidth}{0.4pt}
  \renewcommand{\footrulewidth}{0pt}
}

\fancypagestyle{referencias}{
  \fancyhf{}
  \fancyhead[L]{\includegraphics[height=1.5cm]{portada_img/der.png}}
  \fancyhead[R]{\includegraphics[height=1.5cm]{portada_img/izq.png}}
  \renewcommand{\headrulewidth}{0.4pt}
  \renewcommand{\footrulewidth}{0pt}
}

\newcommand{\var}[1]{\texttt{#1}}
\newcommand{\colab}{\href{https://colab.research.google.com/drive/1uTetGGQgdX5gBBdlnDIoX7KXgtDGWbrd}{Practica 2 en Google Colab}}


\begin{document}
\begingroup
\onehalfspacing

\thispagestyle{empty}
\begin{center}
    \begin{tabular}{m{0.15\textwidth} m{0.60\textwidth} m{0.15\textwidth}}
        \includegraphics[width=2cm]{portada_img/der.png} &
        \centering\Huge\textbf{Universidad Nacional Autónoma de México} &
        \raggedleft\includegraphics[width=2cm]{portada_img/izq.png}
    \end{tabular}
\end{center}

\vspace{0.3cm}
\begin{center}
    {\LARGE \textbf{Facultad de Ingeniería}}\\[1cm]

    {\LARGE \textbf{Alumno:}}\\[0.3cm]
    {\Large Lara Hernandez Angel Husiel - 320060829} \\[1cm]
    {\Large Lugo Manzano Rodrigo - 320206968} \\[1cm]
    {\Large García Cortés Adolfo de Jesús - 320317264} \\[1cm]

    {\LARGE \textbf{Mineria de datos.}} \\[1cm]

	{\LARGE Grupo: 01 - Semestre: 2027-1} \\[1cm]
    
    {\LARGE \textbf{Practica 3}}\\[0.5cm]

    {\LARGE {Machine Learning II}}\\[1cm]

    {\LARGE \textbf{Profesor:}}\\[0.2cm]
    {\Large MTRO. GERARDO GABRIEL CARRASCO ZUÑIGA}\\[1cm]

    {\Large \textbf{Fecha de Entrega:}}\\[0.2cm]
    {\Large 12 de Octubre del 2026}
\end{center}

\endgroup

\clearpage
{
  \makeatletter
  \let\ps@plain\ps@indice
  \makeatother
  \pagestyle{indice}
  \thispagestyle{indice}
  \hypersetup{linkcolor=black}\tableofcontents
  \thispagestyle{indice}
}
\clearpage


\setcounter{page}{1}
\setstretch{1.08}


\section{Objetivo}

Implementar y comparar los algoritmos de minería de patrones frecuentes Apriori y FP-Growth sobre un
conjunto de datos transaccionales de compras, con el fin de extraer, evaluar y optimizar reglas de asociación
que fundamenten una estrategia comercial de ventas cruzadas (cross-selling) y promociones de descuento
mediante criterios multiobjetivo (support, confidence y lift).

\newpage

\section{Carga, exploración e ingeniería de datos transaccionales}

\subsection{Preguntas}

\textbf{¿Cuántas transacciones y cuántos productos únicos componen el catálogo analizado?}

\textbf{¿Cuál es el producto con mayor frecuencia de compra individual en la tienda y cuál es su soporte
empírico?}

\newpage

\section{Extracción de patrones frecuentes: Apriori vs. FP-Growth}

\subsection{Preguntas}

\textbf{¿Coinciden en cantidad y soporte los itemsets frecuentes extraídos por Apriori y FP-Growth?
Justifica teóricamente el resultado.}

\textbf{Explica la diferencia técnica fundamental entre cómo opera Apriori (estrategia por niveles y poda
de candidatos) frente a FP-Growth (árbol comprimido y proyección recursiva). ¿Cuál esperarías
que escale mejor si el catálogo crece a miles de productos?}

\newpage

\section{Generación y filtrado de reglas de asociación}

\subsection{Preguntas}

\textbf{Si una regla presenta una confianza del 88\% pero su Lift es de 1.01, ¿por qué no representa una
sinergia real de compra? Explica la relación entre la probabilidad condicionada P(B|A) y la
probabilidad marginal p(B).}

\textbf{¿Qué diferencia conceptual existe entre una regla con  Lift mayor que 1, Lift igual a 1 y Lift menor que 1?}

\newpage

\section{Visualización del espacio de reglas}

\subsection{Preguntas}

\textbf{A partir del gráfico generado, ¿se concentran las reglas de mayor Lift en la zona de soporte alto o
en la zona de soporte moderado/bajo? Interpreta el fenómeno observado.}

\newpage
\section{Optimización multiobjetivo y propuesta de campaña comercial}

\subsection{Preguntas}

\textbf{¿Por qué seleccionar el top 5 basándose únicamente en el Lift máximo puede resultar riesgoso
para el negocio en un entorno real de retail?}

\textbf{Para la regla clasificada en la posición \#1 de tu ranking, redacta la propuesta comercial explícita
indicando:}

\begin{enumerate}
  \item ¿Cuál producto debe cobrarse a precio regular como ancla (antecedente)?
  \item ¿A cuál producto debe aplicarse el descuento (consecuente)?
  \item ¿Cuál es el alcance esperado (soporte) y la tasa de conversión estimada (confianza) de la
oferta?
\end{enumerate}

\newpage
\section{Conclusión}


\clearpage
\addcontentsline{toc}{section}{Bibliografías}
\begingroup
\makeatletter
\let\ps@plain\ps@referencias
\makeatother
\pagestyle{referencias}
\thispagestyle{referencias}
\renewcommand{\refname}{Bibliografías}
\nocite{referencia}
\bibliographystyle{apacite}
\bibliography{referencias}
\clearpage


\endgroup
\end{document}
